\documentclass[a4paper,11pt]{article}

\usepackage[T1]{fontenc}
\usepackage[utf8]{inputenc}
\usepackage{geometry}
\usepackage{amsmath,amssymb}
\usepackage{xcolor}
\usepackage[markup=underlined,authormarkup=none]{changes}
\usepackage{xspace}

\geometry{margin=2.5cm}
\newenvironment{ReviewerComment}{\begin{quote}\small}{\end{quote}}

\newcommand{\Kinc}{k_{\rm inc}}
\newcommand{\Jinc}{j_{\text{inc}}}
\newcommand{\JcZero}{J_c^0}
\newcommand{\Jeff}{J_{x,c}^{\rm eff}}
\definechangesauthor[name={Response to Reviewers}]{REVISION}

\title{Response to Reviewers}
\author{Alexander Schl\"uter, Ronjit Medda, Ralf M\"uller}
\date{}

\begin{document}

\maketitle

\setlength{\parindent}{0pt}
\setlength{\parskip}{0.55\baselineskip}

We thank both reviewers for their careful reading of the manuscript
and for the constructive suggestions. We have revised the manuscript
accordingly. The main changes are summarized below and are shown in
the highlighted manuscript version.

\section*{Summary of Changes}

\begin{itemize}
  \item Corrected the notation for the velocity of the surfing
  boundary condition from \(\dot{x}_{\rm bc}\) to
  \(\dot{x}^{\rm bc}\) in the text above Eq.~(2).
  \item Added missing definitions of \(J\), \(\psi_{\rm el}\),
  \(\boldsymbol{1}\), \(\boldsymbol{\sigma}\),
  \(\boldsymbol{\varepsilon}^e\),
  \(\boldsymbol{\varepsilon}^p\), and \(M\).
  \item Specified that the norm in the von Mises yield function is
  the Frobenius norm.
  \item Expanded the discussion of the boundary J-integral and
  clarified that it is used as an operational macroscopic energy-flux
  measure for the prescribed numerical experiment, not as a strictly
  path-independent local crack-tip quantity in the classical LEFM
  sense.
  \item Clarified in the figure captions and setup description that
  the model is not a complete compact-tension specimen with locally
  resolved microstructure, but a resolved microstructural region
  embedded in a homogeneous domain and loaded by a prescribed
  near-tip field.
  \item Added that matrix--inclusion interface properties are not
  modeled separately; the interfaces are assumed perfectly bonded.
  \item Added the J-integral contour information to the description
  of Fig.~2 and identified the outer boundary of the plotted domain
  as the contour used for evaluation.
  \item Added a short comment on the absence of crack deflection in
  the present parameter range and related this to the moderate
  inclusion-to-matrix contrasts considered.
\end{itemize}

\section*{Reviewer 1}

\textbf{Reviewer comment.}
\begin{ReviewerComment}
The proceeding ``Effective crack resistance of a simplified ductile
heterogeneous material with circular inclusions'' by Schlüter et al.
provides interesting insights into the influence of microscopic
inclusions on macroscopic crack resistance through numerical
simulations employing homogenization and a phase field model. In
light of the thoroughness of the performed study and the unique
interplay of the inclusion properties it reveals, the proceeding is
suitable for publication. Nevertheless, a few minor queries should be
addressed:
\end{ReviewerComment}

\textbf{Comment 1.}
\begin{ReviewerComment}
A change of \(\dot{x}_{bc}\) to \(\dot{x}^{bc}\) above Eq. (2) would
be consistent with Fig. 1.
\end{ReviewerComment}

\emph{Response.} We agree and have corrected the notation in the
text above Eq.~(2).

\emph{Related change.} Above Eq.~(2):
\[
\replaced{\dot{x}_{\rm bc}}{\dot{x}^{\rm bc}} .
\]
The same notation is also used in the definition of the phase-field
mobility.

\textbf{Comment 2.}
\begin{ReviewerComment}
Some symbols are not properly introduced in the text. These include:
1) \(J\) of Eq. (1),
2) \(\Phi_{\rm el}\), \(\boldsymbol{1}\) and \(\Sigma\) of Eq. (5),
3) \(\epsilon^e\) and \(\epsilon^p\) of Eq. (7),
4) \(M\) of Eq. (15).
\end{ReviewerComment}

\emph{Response.} We agree that several symbols were introduced too
briefly. We have added explicit definitions at the first relevant
locations. In the manuscript notation, the elastic energy density and
stress tensor are denoted by \(\psi_{\rm el}\) and
\(\boldsymbol{\sigma}\), respectively.

\emph{Related changes.}
\begin{itemize}
  \item Introduction, Eq.~(1): \added{Here, \(J\) denotes the
  J-Integral as a measure of the macroscopically observable energetic
  crack-driving force.}
  \item Section~2, after the J-integral: \added{Here,
  \(\psi_{\rm el}\) denotes the degraded elastic energy density
  specified in Section~3.1, \(\boldsymbol{1}\) is the second-order
  identity tensor, and \(\boldsymbol{\sigma}\) is the Cauchy stress
  tensor.}
  \item Section~3.1, after the additive strain decomposition:
  \added{Here, \(\boldsymbol{\varepsilon}^e\) and
  \(\boldsymbol{\varepsilon}^p\) are the elastic and plastic strain
  tensors, respectively.}
  \item Section~3.1, after the phase-field evolution equation:
  \added{The mobility \(M\) controls the relaxation rate of the phase
  field in the pseudo-time evolution equation.}
\end{itemize}

\textbf{Comment 3.}
\begin{ReviewerComment}
It would be helpful for the reader to mention which norm is used in
Eq. (9), for it is not clear from the context.
\end{ReviewerComment}

\emph{Response.} We agree and have specified the norm directly after
the yield function.

\emph{Related change.} Section~3.1:
\[
\added{\|\boldsymbol{s}\|=(\boldsymbol{s}:\boldsymbol{s})^{1/2}.}
\]

\section*{Reviewer 2}

\textbf{Reviewer comment.}
\begin{ReviewerComment}
Overall, I found the manuscript interesting and enjoyable to read.
The proposed methodology is consistent with the framework presented
in Refs. [5,9], extending it to elastoplastic deformation while
investigating the influence of additional parameters. The manuscript
is well structured and presents most of the relevant information in a
clear manner.

For future work, it would be interesting to include a comparison
between a purely elastic and an elastoplastic example in order to
better highlight the additional insights provided by the present
formulation with respect to the original studies.

One question, however, remained unclear to me: What is the
physical/mechanical interpretation of the quantity represented by the
contour integral used in this work?

From the perspective of classical LEFM, the J-integral is a local
crack-tip quantity associated with a single crack tip. For sharp
cracks, path independence requires special treatment of the crack
faces. In the case of diffuse cracks, where no physical crack faces
exist, a simple rectangular contour appears to be sufficient.
However, material heterogeneities and plastic deformation are also
known to influence the interpretation and path independence of the
J-integral. In the literature, additional treatments are often
discussed for these cases. Hosain et al. also appear to acknowledge
this issue (p. 17), although it is not discussed in detail.

Initially, I assumed that the contour integral in the present work
was intended to represent the total energy dissipation rate within
the enclosed region during crack growth rather than a crack-tip
quantity. However, in Hosain et al. (Fig. 4), the effective \(J\) is
validated against Gao's analytical solution for the stress intensity
factor together with Irwin's relation for the crack-tip J-integral.

Therefore, I would appreciate a brief discussion when introducing the
J-integral, together with appropriate references addressing its
application to diffuse cracks and elastoplastic materials. In
particular, it would be helpful to explain why the adopted contour
integral provides a suitable definition of \((J^{\mathrm{eff}})\) in
the present framework, despite the apparent loss of strict path
independence.
\end{ReviewerComment}

\emph{Response.} We thank the reviewer for this positive assessment.
We agree that a direct comparison between purely elastic and
elastoplastic examples would be valuable; however, such a comparison
is beyond the scope of this short proceedings contribution and is a
natural extension for future work. We also agree that the manuscript
should more clearly
state the interpretation of the contour integral. We have revised the
text after the J-integral definition. The revised manuscript now
explains that the boundary integral is not used as a strictly
path-independent local crack-tip quantity in the classical LEFM
sense, particularly when different paths are compared that intersect
the microstructure-resolved region with complex crack patterns and
inelastic material behavior. Instead, following Hossain et al., it is evaluated on the
outer boundary of the prescribed numerical experiment and interpreted
as the macroscopic energy flux required to drive separation of the
resolved microstructure. We also explicitly state that material
heterogeneity, plastic dissipation, and the possible presence of
multiple diffuse cracks may affect path independence within the
microstructure-resolved region. However, path independence is
expected to be restored for contours that enclose the complete
microstructural crack pattern and lie entirely in the idealized
homogeneous embedding domain of effective stiffness, which is
expected to deform elastically. Therefore, the maximum recorded
boundary value is used as a suitable operational definition of the
macroscopic driving force necessary to enforce macroscopic crack
propagation, as it could be observed in an experiment, in the given
microstructure with the given material behavior. It is therefore
interpreted as the effective crack resistance for this specific
numerical experiment.

\emph{Related change.} Section~2, after Eq.~(5):
\begin{quote}
\added{In the present diffuse-crack and heterogeneous elastoplastic
setting, the boundary integral is not used as a strictly
path-independent local crack-tip quantity in the classical LEFM
sense, particularly when different paths are compared that intersect
the microstructure-resolved region with complex crack patterns and
inelastic material behavior. Instead, following Hossain et al., it is evaluated on the
outer boundary of the numerical experiment and interpreted as the
macroscopic energy flux required to drive separation of the enclosed
resolved microstructure. Material heterogeneity, plastic
dissipation, and the possible presence of multiple diffuse cracks may
affect path independence within the microstructure-resolved region.
However, path independence is expected to be restored for contours
that enclose the complete microstructural crack pattern and lie
entirely in the idealized homogeneous embedding domain of effective
stiffness, which is expected to deform elastically. Therefore, the
maximum recorded boundary value is used here as a suitable
operational definition of the macroscopic driving force necessary to
enforce macroscopic crack propagation, as it could be observed in an
experiment, in the given microstructure with the given material
behavior. It is therefore interpreted as the effective crack
resistance for the prescribed numerical experiment.}
\end{quote}

\textbf{Additional suggestions}

\textbf{Additional suggestion 1.}
\begin{ReviewerComment}
1.In Fig. 2, it would be helpful to state more explicitly that the
model does not represent a complete CT specimen containing a locally
resolved substructure. Instead, only the microstructure within a
homogeneous domain is modeled explicitly, while the near-tip field
corresponding to a crack with constant stress intensity factor is
imposed through the boundary conditions. Although this is explained
in the manuscript, the figure may initially suggest a different
interpretation.
\end{ReviewerComment}

\emph{Response.} We agree and have clarified this point in the
captions of both setup figures. The revised text now states that the
figures show a resolved microstructural region embedded in a
homogeneous domain and subjected to a prescribed near-tip field, not
a complete compact-tension specimen.

\emph{Related changes.}
\begin{itemize}
  \item Caption of Fig.~1: \added{The sketch illustrates the imposed
  near-tip boundary field around the resolved region and should not
  be read as a complete compact-tension specimen with locally
  resolved substructure.}
  \item Caption of Fig.~2: \added{The model resolves the
  microstructure inside a homogeneous embedding domain subjected to a
  prescribed near-tip field, not a complete compact-tension
  specimen.}
\end{itemize}

\textbf{Additional suggestion 2.}
\begin{ReviewerComment}
2.Similar to the underlying references, it would be useful to
explicitly mention that interface properties between the inclusions
and the matrix are neglected.
\end{ReviewerComment}

\emph{Response.} We agree and have added this assumption to the
geometry and material-parameter section.

\emph{Related change.} Section~4.1:
\begin{quote}
\added{The interfaces between inclusions and matrix are assumed
perfectly bonded; separate interface stiffnesses, strengths, or
fracture resistances are not introduced.}
\end{quote}

\textbf{Additional suggestion 3.}
\begin{ReviewerComment}
3.Fig. 2 would benefit from showing the contour used for the
J-integral evaluation together with the principal dimensions of the
model (e.g., \((L)\), \((w_s)\), ligament width, etc.). This would
make the numerical setup easier to understand.
\end{ReviewerComment}

\emph{Response.} We have clarified the contour used for the
J-integral evaluation in the figure caption and accompanying setup
description. In the present formulation, \(J_x\) is evaluated on the
outer boundary of the plotted computational domain. The principal
geometric dimensions are already defined in the text: inclusion
diameter \(L\), ligament width \(w_s\), cell size
\(w_c=L+w_s\), and the \(6\times3\) resolved inclusion array. We
therefore chose a caption/text clarification rather than redrawing
the figure for this minor revision.

\textbf{Additional suggestion 4.}
\begin{ReviewerComment}
4.I initially expected crack deflection to occur for stiff and tough
inclusions, particularly since the considered configuration appears
comparable to the offset configuration shown in Fig. 11 of Hosain et
al. [5]. It would therefore be helpful to comment on whether the
absence of crack deflection is primarily a consequence of the
relatively moderate inclusion-to-matrix property contrast considered
in the present study. Alternatively, a brief reference to the
corresponding observations in Ref. [5] would also help readers
interpret the results.
\end{ReviewerComment}

\emph{Response.} We agree and have added a short comment before the
geometry subsection. The revised manuscript states that no crack
deflection around stiff and tough inclusions was observed in the
present parameter range, and that this is consistent with the
moderate inclusion-to-matrix contrasts considered. We also note that
stronger contrasts or different offset configurations, such as those
discussed by Hossain et al., may lead to different crack-path
selection.

\emph{Related change.} Section~4:
\begin{quote}
\added{Crack deflection around stiff and tough inclusions was not
observed in the present parameter range. This is consistent with
interpreting the investigated inclusion-to-matrix property contrasts
as moderate; stronger contrasts or different offset configurations,
such as those discussed by Hossain et al., may lead to different
crack-path selection.}
\end{quote}

\end{document}
